\labsection{2}{Arrays, multidimensional storage, and matrices}{sec:lab02}

\begin{sourcealignment}{Official Lab Manual --- Lab II}
\textbf{Objective.} Understand the array as a basic prerequisite for more complex data structures.

\textbf{Outcomes.} Explain how arrays are stored in memory and use arrays to solve programming problems.

\textbf{Problem directions.} Real-world scenarios involving arrays; multidimensional arrays; matrix-related problems such as matrix multiplication.
\end{sourcealignment}

\begin{labproject}{Process and multiply matrices}
\textbf{Coverage.} Chapter 2.\\
\textbf{Goal.} Multiply two matrices and explain what each of the three nested loops controls.

\begin{lstlisting}
#include <iostream>
using namespace std;

int main() {
    const int R = 2, K = 3, C = 2;
    int A[R][K] = {{1,2,3}, {4,5,6}};
    int B[K][C] = {{7,8}, {9,10}, {11,12}};
    int P[R][C] = {};

    for (int i = 0; i < R; ++i)
        for (int j = 0; j < C; ++j)
            for (int t = 0; t < K; ++t)
                P[i][j] += A[i][t] * B[t][j];
}
\end{lstlisting}

\begin{tryit}
Add two of the following: row sums, column sums, transpose, largest element, diagonal sum, or a check for whether a square matrix is symmetric.
\end{tryit}

\begin{workedexample}
For matrix multiplication, the result size is $R\times C$. Each output cell $P[i][j]$ accumulates products across the shared dimension $K$.
\end{workedexample}
\end{labproject}

\begin{verification}
\begin{itemize}
  \item Validate that the number of columns in $A$ equals the number of rows in $B$.
  \item Print both inputs and the output matrix with clear row/column formatting.
  \item Trace one output cell by hand and match it to the program result.
  \item Explain what each nested loop controls.
\end{itemize}
\end{verification}

\begin{labdeliverable}
Submit source code, output, and one hand-worked multiplication cell showing the same arithmetic performed by the innermost loop.
\end{labdeliverable}
